\section{Data-driven \texorpdfstring{$\ttbar$+tW}{ttbar+tW} estimate from the \texorpdfstring{$e\mu$}{e-mu} sideband}
\label{app:emu_method}

This appendix repeats, for 2024, the study of App.~J of Ref.~\cite{bib:AN19083}:
the $\ttbar$+tW yield of a same-flavor region is estimated from the data of the
$e\mu$ sideband with the same kinematics,
%
\begin{equation}
  N_{\ttbar}(\text{same flavor}) =
  \bigl[N_{\text{data}}(e\mu) - N_{\text{other bkg}}^{\text{MC}}(e\mu)\bigr]
  \times R_{ee,\mu\mu/e\mu},
  \label{eq:emu_method}
\end{equation}
%
with the flavor ratio $R$ taken from the $\ttbar$+tW simulation. It is the
method of the earlier analysis and serves here as a cross-check of the
simultaneous fit of Sec.~\ref{subsec:top_draft}, which is the method used for
the results; neither the ratios nor the closure uncertainties derived below
enter the fit.

The cut-based analyzer output contains no region below $\mlljj$ ($\mlJ$)
$= 800\GeV$, so all regions of this appendix are rebuilt from the event-level
output of the BDT preselection, which applies the same object and event
selection without the mass requirements: $\mll > 400\GeV$ ($200\GeV$) and
$\mlljj$ ($\mlJ$) $> 800\GeV$ define the resolved (boosted) high-mass regions,
and $\mll > 200\GeV$ with $\mlljj$ ($\mlJ$) $< 800\GeV$ the low-mass control
regions. In the resolved high-mass $e\mu$ region the $\ttbar$ simulation
reproduces the cut-based histograms to $0.1\%$ in the number of selected
events. The same-flavor data are used in the low-mass regions only; the
signal regions remain blinded.

\subsection{Signal contamination}
The derivation assumes that the $e\mu$ sidebands are free of signal.
Table~\ref{tab:emu_sigcont} gives the simulated signal yield relative to the
total background for the mass points available in the preselection output.
The resolved sideband and the sideband whose in-jet lepton differs in flavor
from the signal leptons contain less than $0.1\%$ of signal for
$\mWR \geq 3\TeV$. The $ee$ ($\mu\mu$) signal does populate the
$e$+($\mu$-jet) ($\mu$+($e$-jet)) sideband, through a nonprompt lepton of
the other flavor inside the \N{}-candidate jet: $0.6\%$ ($0.05\%$) at
$(3000, 200)\GeV$ and $3.5\%$ ($1.0\%$) at $(2000, 200)\GeV$; the fractions
at $\mWR = 1\TeV$, a mass excluded by the earlier searches, reach $36\%$.

\begin{table}[htbp]
  \centering
  \topcaption{Signal yield in the high-mass $e\mu$ sidebands in percent of the
    total simulated background, 2024, for the $ee$ (top) and $\mu\mu$ (bottom)
    signal.}
  \label{tab:emu_sigcont}
  \small
  \input{tables/05_backgrounds/2024/signal_contamination_2024}
\end{table}

\subsection{Flavor ratios}
Figures~\ref{fig:emu_ratio_resolved}--\ref{fig:emu_ratio_boosted_ejet} show the
ratio of the $\ttbar$+tW simulation in the same-flavor region to the $e\mu$
sideband as a function of the reconstructed \PWR{} mass, in the low- and the
high-mass regions, with a fit of a constant. The boosted $\mu\mu$ region is
compared with the $e$+($\mu$-jet) sideband and the boosted $ee$ region with the
$\mu$+($e$-jet) sideband, so that the lepton inside the jet has the same flavor
in the numerator and the denominator. The fitted constants are collected in
Table~\ref{tab:emu_ratio_fit}. The low- and high-mass values agree within
$5\%$; the resolved $\mu\mu$ ratio and the low-mass boosted ratios are not
well described by a constant ($\chi^{2}/\text{ndf}$ of $27/5$ to $29/5$), the
ratio rising by a few percent with the mass.

\begin{table}[htbp]
  \centering
  \topcaption{Constant fits of the same-flavor to $e\mu$ ratio of the
    $\ttbar$+tW simulation, 2024, with statistical uncertainties.}
  \label{tab:emu_ratio_fit}
  \input{tables/05_backgrounds/2024/emu_ratio_fit_2024}
\end{table}

\begin{figure}[htbp]
  \centering
  \includegraphics[width=0.45\textwidth]{fig/05_backgrounds/2024/emu_method/ratio_resolved_low.png}
  \includegraphics[width=0.45\textwidth]{fig/05_backgrounds/2024/emu_method/ratio_resolved_high.png}
  \caption{Signal-region to $e\mu$-sideband ratio of the $\ttbar$+tW simulation
    as a function of $\mlljj$ in the resolved category, 2024, in the low-mass
    (left) and high-mass (right) regions.}
  \label{fig:emu_ratio_resolved}
\end{figure}

\begin{figure}[htbp]
  \centering
  \includegraphics[width=0.45\textwidth]{fig/05_backgrounds/2024/emu_method/ratio_boosted_mujet_low.png}
  \includegraphics[width=0.45\textwidth]{fig/05_backgrounds/2024/emu_method/ratio_boosted_mujet_high.png}
  \caption{Same as Fig.~\ref{fig:emu_ratio_resolved} as a function of $\mlJ$ in
    the boosted category with a $\mu$-jet ($\mu\mu$ over $e$+($\mu$-jet)).}
  \label{fig:emu_ratio_boosted_mujet}
\end{figure}

\begin{figure}[htbp]
  \centering
  \includegraphics[width=0.45\textwidth]{fig/05_backgrounds/2024/emu_method/ratio_boosted_ejet_low.png}
  \includegraphics[width=0.45\textwidth]{fig/05_backgrounds/2024/emu_method/ratio_boosted_ejet_high.png}
  \caption{Same as Fig.~\ref{fig:emu_ratio_resolved} in the boosted category
    with an $e$-jet ($ee$ over $\mu$+($e$-jet)).}
  \label{fig:emu_ratio_boosted_ejet}
\end{figure}

\subsection{\texorpdfstring{$e\mu$}{e-mu} sidebands in data}
Figure~\ref{fig:emu_sidebands} compares data and simulation in the three
$e\mu$ sidebands, in the low- and the high-mass regions. The band contains the
statistical and the weight-based systematic uncertainties (lepton, trigger,
pileup and $\LSFiii$ scale factors); the jet and lepton energy scale and
resolution variations are not available in this output.

\begin{figure}[htbp]
  \centering
  \includegraphics[width=0.40\textwidth]{fig/05_backgrounds/2024/emu_method/Res_EM_low.png}
  \includegraphics[width=0.40\textwidth]{fig/05_backgrounds/2024/emu_method/Res_EM_high.png}\\
  \includegraphics[width=0.40\textwidth]{fig/05_backgrounds/2024/emu_method/Boost_EMJ_low.png}
  \includegraphics[width=0.40\textwidth]{fig/05_backgrounds/2024/emu_method/Boost_EMJ_high.png}\\
  \includegraphics[width=0.40\textwidth]{fig/05_backgrounds/2024/emu_method/Boost_MEJ_low.png}
  \includegraphics[width=0.40\textwidth]{fig/05_backgrounds/2024/emu_method/Boost_MEJ_high.png}
  \caption{Reconstructed \PWR{} mass in the flavor sidebands, 2024: resolved
    $e\mu$ (top), boosted $e$+($\mu$-jet) (middle) and boosted $\mu$+($e$-jet)
    (bottom), in the low-mass (left) and high-mass (right) regions.}
  \label{fig:emu_sidebands}
\end{figure}

\subsection{Closure in the low-mass regions}
The method is tested in the low-mass same-flavor regions, where the
$\ttbar$+tW prediction of Eq.~(\ref{eq:emu_method}), with the low-mass ratio of
Table~\ref{tab:emu_ratio_fit}, is added to the simulated other backgrounds and
compared with data (Fig.~\ref{fig:emu_lowmass_dd}).
In the resolved regions, where the data-driven $\ttbar$+tW makes up $80\%$ of
the prediction, the observed-to-predicted ratio is $1.02$ ($ee$) and $1.08$
($\mu\mu$), and between $0.94$ and $1.11$ in every bin above $300\GeV$. The
data-driven $\ttbar$+tW yield is $0.91$ of the simulated one. In the boosted
regions the ratio is $1.52$ ($ee$) and $1.45$ ($\mu\mu$). There the
data-driven $\ttbar$+tW is only $30\%$ ($38\%$) of the prediction, the rest
being DY and nonprompt simulation, so the disagreement cannot be attributed to
the $e\mu$ method: even a $60\%$ uncertainty on the data-driven component moves
the prediction by $20\%$. \TODO{The boosted low-mass regions point to the DY and
nonprompt modeling at low $\mlJ$, in line with the boosted DY CR2 closure of
Sec.~\ref{subsubsec:dyvalidation_draft}; follow up before using these regions
to assign an uncertainty.}

Following the Run~2 study, Figs.~\ref{fig:emu_obs_pred_res_ee}--\ref{fig:emu_obs_pred_boost_mm}
show the observed-to-predicted ratio when a systematic uncertainty between
$5\%$ and $60\%$ is assigned to the data-driven $\ttbar$+tW component. In
the resolved regions $15\%$ covers all bins above $300\GeV$ (Run~2:
$20\%$); in the boosted regions no value in the scanned range covers the
discrepancy, for the reason given above (Run~2: $30\%$).

\begin{figure}[htbp]
  \centering
  \includegraphics[width=0.40\textwidth]{fig/05_backgrounds/2024/emu_method/lowmass_dd_Res_EE.png}
  \includegraphics[width=0.40\textwidth]{fig/05_backgrounds/2024/emu_method/lowmass_dd_Res_MM.png}\\
  \includegraphics[width=0.40\textwidth]{fig/05_backgrounds/2024/emu_method/lowmass_dd_Boost_EE.png}
  \includegraphics[width=0.40\textwidth]{fig/05_backgrounds/2024/emu_method/lowmass_dd_Boost_MM.png}
  \caption{Reconstructed \PWR{} mass in the resolved (top) and boosted (bottom)
    low-mass control regions for $ee$ (left) and $\mu\mu$ (right) events, 2024,
    with the $\ttbar$+tW contribution estimated from the $e\mu$ sideband data.
    The band contains the statistical uncertainty of the prediction and the
    weight-based systematic uncertainties of the simulated backgrounds.}
  \label{fig:emu_lowmass_dd}
\end{figure}

\begin{figure}[htbp]
  \centering
  \includegraphics[width=\textwidth]{fig/05_backgrounds/2024/emu_method/obs_pred_Res_EE.png}
  \caption{Observed-to-predicted ratio in the resolved low-mass $ee$ region
    versus $\mlljj$, 2024, for a systematic uncertainty of $5\%$ to $60\%$ on
    the data-driven $\ttbar$+tW estimate. Orange: statistical uncertainty of the
    prediction; gray: statistical and systematic.}
  \label{fig:emu_obs_pred_res_ee}
\end{figure}

\begin{figure}[htbp]
  \centering
  \includegraphics[width=\textwidth]{fig/05_backgrounds/2024/emu_method/obs_pred_Res_MM.png}
  \caption{Same as Fig.~\ref{fig:emu_obs_pred_res_ee} for the resolved low-mass
    $\mu\mu$ region.}
  \label{fig:emu_obs_pred_res_mm}
\end{figure}

\begin{figure}[htbp]
  \centering
  \includegraphics[width=\textwidth]{fig/05_backgrounds/2024/emu_method/obs_pred_Boost_EE.png}
  \caption{Same as Fig.~\ref{fig:emu_obs_pred_res_ee} for the boosted low-mass
    $ee$ region, versus $\mlJ$.}
  \label{fig:emu_obs_pred_boost_ee}
\end{figure}

\begin{figure}[htbp]
  \centering
  \includegraphics[width=\textwidth]{fig/05_backgrounds/2024/emu_method/obs_pred_Boost_MM.png}
  \caption{Same as Fig.~\ref{fig:emu_obs_pred_res_ee} for the boosted low-mass
    $\mu\mu$ region, versus $\mlJ$.}
  \label{fig:emu_obs_pred_boost_mm}
\end{figure}
